% Lösungen Fokus Nullstellen – Einheit 4, Hauptnummern 6 bis 18
\einheitenkopf{Einheit 4 von 4 · Nullstellen und Schnittpunkte berechnen}
\erg{6}{a)~$(x - 5)^2 - 4 = 0$ \quad b)~$x^2 + 7x - 8 = 0$ \quad c)~$-2x^2 + 8 = 0$ \quad d)~$x^2 - 1 = x + 5$ \quad e)~$(x + 1)^2 = -x + 3$}
\erg{7}{a)~$x_1 = 5$, $x_2 = -5$ \quad b)~$x - 4 = \pm 1$; $x_1 = 5$, $x_2 = 3$ \quad c)~$x - 6 = \pm 2$; $x_1 = 8$, $x_2 = 4$ \quad d)~$x - 3 = \pm 5$; $x_1 = 8$, $x_2 = -2$}
\erg{8}{a)~$x + 2 = \pm 3$; $x_1 = 1$, $x_2 = -5$ \quad b)~$x + 5 = \pm 2$; $x_1 = -3$, $x_2 = -7$ \quad c)~$(x - 7)^2 = 0$; nur $x = 7$ \quad d)~$(x + 1)^2 = -4$; keine Nullstelle}
\erg{9}{a) zwei \quad b) keine \quad c) eine \quad d) zwei \quad e) keine \quad f) Nullstellen liegen auf der x-Achse, und dort ist der y-Wert $f(x) = 0$. \quad g) Nein: Die Parabel ist nach unten geöffnet, ihr Scheitel $S(4 \mid -1)$ ist ihr höchster Punkt und liegt schon unter der x-Achse – sie hat keine Nullstelle.}
\erg{10}{a)~$x_{1,2} = 1 \pm 4$; $x_1 = 5$, $x_2 = -3$ \quad b)~$x_{1,2} = 4 \pm \sqrt{16 - 12} = 4 \pm 2$; $x_1 = 6$, $x_2 = 2$ \quad c)~$x_{1,2} = -2 \pm \sqrt{4 + 21} = -2 \pm 5$; $x_1 = 3$, $x_2 = -7$ \quad d)~$x_{1,2} = -5 \pm \sqrt{25 - 16} = -5 \pm 3$; $x_1 = -2$, $x_2 = -8$}
\erg{11}{a)~$x_{1,2} = 2{,}5 \pm \sqrt{6{,}25 - 6} = 2{,}5 \pm 0{,}5$; $x_1 = 3$, $x_2 = 2$ \quad b)~$x_{1,2} = -1{,}5 \pm \sqrt{2{,}25 + 10} = -1{,}5 \pm 3{,}5$; $x_1 = 2$, $x_2 = -5$ \quad c)~$x_{1,2} = 5 \pm \sqrt{25 - 25} = 5$; nur eine Nullstelle $x = 5$ \quad d)~$x_{1,2} = -2 \pm \sqrt{4 - 7} = -2 \pm \sqrt{-3}$; keine Nullstelle}
\erg{12}{a)~$x_{1,2} = 1 \pm \sqrt{3}$; $x_1 \approx 2{,}73$, $x_2 \approx -0{,}73$ \quad b)~$x_{1,2} = -3 \pm \sqrt{7}$; $x_1 \approx -0{,}35$, $x_2 \approx -5{,}65$ \quad c)~$x_{1,2} = 2 \pm \sqrt{4 + 2} = 2 \pm \sqrt{6}$; $x_1 \approx 4{,}45$, $x_2 \approx -0{,}45$}
\erg{13}{a)~$x_{1,2} = -1 \pm 3$; $x_1 = 2$, $x_2 = -4$ \quad b) durch 3: $x^2 - 7x + 10 = 0$; $x_{1,2} = 3{,}5 \pm \sqrt{12{,}25 - 10} = 3{,}5 \pm 1{,}5$; $x_1 = 5$, $x_2 = 2$ \quad c) durch $-1$: $x^2 - 4x - 5 = 0$; $x_{1,2} = 2 \pm \sqrt{4 + 5} = 2 \pm 3$; $x_1 = 5$, $x_2 = -1$}
\erg{14}{a)~$2(x - 3)^2 = 18$, $(x - 3)^2 = 9$, $x - 3 = \pm 3$; $x_1 = 6$, $x_2 = 0$ \quad b)~$-(x + 4)^2 = -1$, $(x + 4)^2 = 1$, $x + 4 = \pm 1$; $x_1 = -3$, $x_2 = -5$}
\erg{15}{a) durch 2: $x^2 - x - 12 = 0$; $x_{1,2} = 0{,}5 \pm \sqrt{0{,}25 + 12} = 0{,}5 \pm 3{,}5$; $x_1 = 4$, $x_2 = -3$ \quad b)~$x_{1,2} = -2 \pm \sqrt{4 + 1} = -2 \pm \sqrt{5}$; $x_1 \approx 0{,}24$, $x_2 \approx -4{,}24$ \quad c)~$(x - 4)^2 - 7 = x^2 - 8x + 16 - 7 = x^2 - 8x + 9$ (wA); $x_{1,2} = 4 \pm \sqrt{16 - 9} = 4 \pm \sqrt{7}$; $x_1 \approx 6{,}65$, $x_2 \approx 1{,}35$}
\erg{16}{a) Fehler: Beim Wurzelziehen fehlt die zweite Möglichkeit $x + 1 = -7$. Richtig: $x + 1 = 7$ oder $x + 1 = -7$; $x_1 = 6$, $x_2 = -8$ \quad b)~$x - 5 = \pm 6$; $x_1 = 11$, $x_2 = -1$}
\erg{17}{a) Fehler: Ben hat nicht zuerst durch 2 geteilt; die p-q-Formel gilt nur, wenn vor $x^2$ keine Zahl steht. Richtig: $x^2 + 6x + 5 = 0$; $x_{1,2} = -3 \pm \sqrt{9 - 5} = -3 \pm 2$; $x_1 = -1$, $x_2 = -5$ \quad b) durch 3: $x^2 - 4x - 12 = 0$; $x_{1,2} = 2 \pm \sqrt{4 + 12} = 2 \pm 4$; $x_1 = 6$, $x_2 = -2$}
\erg{18}{a)~$h(0) = -0{,}1 \cdot 16 + 3{,}6 = 2$; der Ball verlässt die Hand in 2\,m Höhe. \quad b)~$h(x) = 0$: $(x - 4)^2 = 36$, $x - 4 = \pm 6$; $x_1 = 10$, $x_2 = -2$; der Ball kommt 10\,m von Tim entfernt auf. \quad c)~$x = -2$ läge 2\,m hinter Tim; der Ball fliegt aber erst ab $x = 0$ (aus der Hand) nach vorn, deshalb passt nur $x = 10$.}
